\chapter{TFT o TFG de investigación}
Para trabajos de investigación, especialmente en el contexto de Trabajos de Fin de Título (TFT) o Trabajos de Fin de Grado (TFG), es crucial estructurar el documento de manera que refleje el rigor científico y la profundidad del estudio realizado. Aquí se detallan las secciones sugeridas: 
\begin{enumerate}
    \item Introducción: Esta sección establece el escenario para el estudio, introduciendo el tema de investigación, su relevancia dentro del campo de la informática o ciencia de datos, y la motivación detrás del estudio. Debe captar el interés del lector y proporcionar una visión general clara de los objetivos del trabajo. 
    \item Descripción del problema: Se debe detallar con precisión el problema específico que se aborda, explicando su importancia y cómo se inscribe dentro de los desafíos actuales del área. Es fundamental argumentar la necesidad de una investigación en esta área.
    \item Objetivos: Enumera los objetivos generales y específicos del proyecto. Los objetivos deben ser concisos y medibles, proporcionando una guía clara para las etapas de desarrollo que siguen.
    \item Marco teórico/referencial: Esta sección es esencial para establecer las bases del estudio. Debe incluir una revisión exhaustiva de la literatura existente, teorías, modelos, y trabajos previos que sean relevantes para el tema de investigación. El marco teórico no solo contextualiza el estudio dentro del campo sino que también identifica lagunas en la investigación existente que el trabajo actual busca abordar. 
    \item Metodología y plan de trabajo: En esta parte, se describe detalladamente cómo se lleva a cabo la investigación. Incluye la elección de métodos de investigación, diseño de experimentos, técnicas de recopilación y análisis de datos, y herramientas utilizadas. Esto puede abarcar desde la definición de algoritmos hasta el uso de frameworks específicos para el análisis de datos. 
    \item Propuesta de solución: Aquí se detalla la solución propuesta para el problema de investigación, explicando cómo se desarrolló, los fundamentos teóricos o técnicos en los que se basa, y cómo se diferencia o mejora las soluciones existentes. 
    \item Experimentos y análisis de resultados: Esta sección debe describir los experimentos realizados para validar la propuesta de solución, incluyendo la configuración, los datos utilizados, las métricas de evaluación, y los resultados obtenidos. El análisis debe ser crítico, destacando tanto los éxitos como las limitaciones de la propuesta basándose en los datos recogidos. 
    \item Conclusiones: Finalmente, se resumen las principales contribuciones del trabajo, se reconocen sus limitaciones y se sugieren direcciones para investigaciones futuras.
\end{enumerate}